
% =====================================================================
\section{What to take to MTEC}
\label{sec:discussion}
% =====================================================================

\subsection{What was found}

Ohta's measures survive out of sample. On a year his study did not cover, round prices take 14.95\% of large-trade volume against 10.77\% of small-trade volume; the large-minus-small ordering his mechanism predicts holds (significant at the 1\% level); finer grids cluster more; the paper's size gradient does not appear in this screened universe --- documented with its tick-mix decomposition in Section~\ref{sec:stylized} rather than assumed away; and trades at round prices carry a positive impact premium of 0.43 basis points at one minute --- about 34\% of the paper's published magnitude, and transient where his horizon ends. None of the comparisons was tuned; the benchmarks were fixed before the numbers were computed.

On the question the paper leaves open, clustering carries a measurable association with next-day trading costs beyond the outcome's own one-day lag (0.0188 on the log effective spread, $t=2.51$; depth complements it at -0.0441), and a placebo across the other last digits leaves 0 of eight at conventional significance. The robustness work bounds the claim: the pooled estimate is significant at the 5\% level, but it falls below conventional significance once the outcome's history means five lags rather than one; its sign reverses between the within-stock and cross-sectional designs. The within-day version --- 0.0189 ($t=8.80$) under stock-day and bucket fixed effects --- is the form of the relationship this study would defend without qualification.

The book-based measures are the clearer contribution, read with discipline. Round prices hold 13.21\% of standing visible depth against a uniform 10\%, observing the stale inventory directly; standing round-price depth predicts round-price execution ($t=11.4$), though partly mechanically; and the inferred flow shares are nulls under severe ten-level censoring, which adjudicates nothing. The inferred submission shares land on the paper's levels without calibration (0.102 and 0.103 against his 0.106 and 0.105), which is the strongest external validation in the study; the cancellation ratios overshoot (0.91 and 0.90 against 0.806 and 0.812), a real difference --- more cancellation in 2025, or netting bias in the ten-level window --- that is reported as open rather than absorbed into the validation claim.

The strategy demonstration prices the idea rather than selling it: gross -3.1 basis points a day against costs of 32.0 at daily rebalancing, and smoothing to five days cuts the costs to 13.5 without uncovering alpha. The economics point at execution, not selection.

\subsection{What it is not}

\paragraph{Not causal.}
Four months, no experiment, no instrument. Everything reported is a descriptive
association or a predictive relationship. Clustering and liquidity are plausibly
both driven by things this study does not observe, and nothing here identifies a
direction of causation.

\paragraph{Not Ohta's identification.}
He identifies noise traders with margin balances and ownership shares. Those are
unavailable here, so this study leans entirely on the claim his conclusion
advances --- that the clustering measures are themselves the proxy. That makes
the exercise a test of his proposal rather than an independent confirmation of
it, and a genuine replication would need the ownership data.

\paragraph{A narrow window on the book.}
Ten visible levels sound like a lot and are not. On the stock-days examined in
detail, 81\% of quoted volume sits beyond them. The inferred order flow
describes the neighbourhood of the best quote, which is where execution happens,
but it is silent about the far book --- and the far book is exactly where Ohta
argues the interesting stale orders sit.

\paragraph{A sample shaped by its filters.}
Excluding non-power-of-ten tick sizes removes the 0.5-yen grid, and with it much
of the mid-priced TOPIX 500 population; the 9:10 rule thins the panel on
shock days. Both follow from the paper's filters rather than from choices made
here, but together they mean the regression sample is not the market, and is
least representative exactly when markets are wildest.

\subsection{Mapping the results onto the internship brief}

The brief --- a rebalancing strategy built on this indicator --- has three
defensible readings, and the results here speak to each differently.

\paragraph{(A) Execution-aware rebalancing.} The best-supported use. The
round-price impact premium is positive, front-loaded and transient
(Section~\ref{sec:stylized}), and the within-day clustering--spread link is
the study's most robust regression (Section~\ref{sec:book}). Both are
statements about \emph{what an order pays and when}, which is what an
execution schedule consumes: avoid crossing at round levels, avoid resting own
orders at them, time non-urgent legs away from high-clustering buckets. No
forecasting is required and a basis point is a real saving.

\paragraph{(B) A signal tilt.} The least supported. The daily association is
real but the pooled estimate is significant at the 5\% level, but it falls below conventional significance once the outcome's history means five lags rather than one; its sign reverses between the within-stock and cross-sectional designs,
the cross-sectional sign points the other way (the Fama--MacBeth estimate), and
the costed demonstration in Section~\ref{sec:strategy} shows the daily
version pays its gross away several times over while the smoothed version has
nothing left to pay with. If a tilt is wanted anyway, it needs months of
holding period, size discipline, and an out-of-sample year.

\paragraph{(C) Rebalance-trigger design.} Untested here and testable with
this exact pipeline: modulate rebalancing bands or urgency by the residualised
clustering state, on the argument that rebalancing is liquidity provision and
the counterparty is more likely uninformed when noise traders are active. The
panel and cost machinery this prototype builds is the spine that experiment
needs.

\paragraph{Data to ask for on day one.} Margin-trading balances (the
\emph{nisshoukin} securities-finance series) and ownership shares would
convert the proxy test into a joint test of proxy and mechanism; a second year
of tape would give the sub-period and out-of-sample splits this sample is too
short for; and the full FLEX feed would lift the ten-level censoring that
mutes the order-flow variables.

\subsection{Conclusion}

The sentence at the end of Ohta (2026) proposes that price clustering can serve
as an observable proxy for noise-trader activity, and that its relationship to
price formation and liquidity could then be studied at daily or shorter
frequency. On the January--April 2025 Tokyo Stock Exchange tape the proposal survives its
first contact with data the paper never saw: the measures replicate, the
within-day liquidity relationship is robust, and the standing round-price
inventory the mechanism requires is directly visible in the book. The daily
version of the relationship is thinner --- the pooled estimate is significant at the 5\% level, but it falls below conventional significance once the outcome's history means five lags rather than one; its sign reverses between the within-stock and cross-sectional designs --- and the indicator
does not support a standalone trading rule at this horizon, because the costs
exceed the signal. The most defensible use of the measure is the one that needs
no forecast at all: knowing what a round-price execution costs, and scheduling
around it.
